% Responsable(s): por asignar
% Objetivo: presentar el problema, la motivación, el objetivo de la revisión
% y la pregunta de investigación.

% Redacción pendiente.
El rendimiento de los sistemas computacionales actuales no depende solamente
de la capacidad de procesamiento de la CPU, sino también de la rapidez con la
que la jerarquía de memoria puede suministrar los datos requeridos. A pesar de
los avances en la arquitectura de los procesadores, la latencia de acceso a la
memoria DRAM no ha disminuido significativamente en las últimas décadas, por
lo que continúa representando un posible cuello de botella para el rendimiento
del sistema. En arquitecturas multinúcleo, este problema puede intensificarse
debido a la contención producida cuando varios núcleos intentan acceder de
forma simultánea a bancos de memoria compartidos
\cite{alawneh2023dram}.

Una de las principales consecuencias de esta diferencia de velocidades es la 
aparición de \textit{memory stalls}, durante los cuales el procesador debe 
detener temporalmente la ejecución de intrucciones. Debido a esta relación,
la latencia de acceso no solo constituye una característica propia de la 
memoria, sino que también repercute directamente sobre métricas del rendimiento
del procesador. Park et al. muestran que una gestión coordinada entre la CPU 
y los diferentes niveles de la jerarquía de memoria puede reducir los ciclos 
asociados a estas esperas \cite{park2022dvfs}.

El análisis del rendimiento de una memoria tampoco puede limitarse a su ancho 
de banda teórico. En sistemas basados en DDR4 existen retardos asociados con
operaciones internas como los ciclos de refresco, activación y precarga. Estos
factores pueden reducir el throughput que realmente puede aprovechar una aplicación.
Shahbazi y Tsui demostraron experimentalmente esta diferencia al implementar 
cuatro memorias DDR4-2400, donde esto alcanzo un throughput de 135.52~Gbps y 
mejorar aproximadamente un 40\% el ancho de banda \cite{shahbazi2026ddr4}.

Asimismo, el ancho de banda y la latencia deben estudiarse de manera conjunta.
Esmaili-Dokht et al. caracterizan sistemas reales equipados con tecnologías
DDR4, DDR5, HBM2, Y HBM2E mediante curvas de ancho de banda. Sus resultados
muestran que la latencia de acceso varía conforme aumenta la presión sobre el
sistema de memoria. Además, el rendimiento no depende exclusivamente de la
tecnología DRAM, sino también de elementos como la jerarquía de caché, el 
controlador de memoria y la arquitectura del procesador \cite{esmailidokht2024mess}.
A partir de este contexto, el presente trabajo analiza el tiempo de acceso y el
rendimiento de diferentes tecnologías y organizaciones de memoria, con énfasis
en la relación entre latencia, ancho de banda y rendimiento del sistema.
